\section{Interpretador e JIT}
\label{sec:jit}

A HotSpot começa interpretando o bytecode. Enquanto o programa roda, ela
mede onde o tempo é gasto e compila para código de máquina só o que é
quente; a documentação da Oracle descreve o compilador como
\emph{adaptativo} e diz que a maior parte do código, pouco usada, nem
chega a ser compilada (\citetitle{doc:jvm-overview}). Esse compilador é o
\termo{JIT}{\emph{Just-in-time}: compila durante a execução, com
informação sobre como o programa está de fato rodando.}.

\paragraph{Tiered compilation.}
Há dois compiladores, que o manual do comando \texttt{java}
(\citetitle{doc:java-cmd}) chama de C1 e C2. Com
\termo{tiered compilation}{Compilação em camadas: o método passa do
interpretador para o C1 (rápido de compilar, com coleta de perfil) e,
se continuar quente, para o C2 (mais lento de compilar, código mais
otimizado).}, introduzida no Java SE~7 e ligada por padrão, o
\emph{client compiler} (C1) gera versões que coletam perfil sobre si
mesmas, e o \emph{server compiler} (C2) usa esse perfil para compilar de
forma agressiva (\citetitle{doc:hotspot-perf}).

\paragraph{Interpretador, JIT ou AOT: quem decide, e quanto custa no startup?}
Como a JVM decide entre interpretador, JIT e AOT, e qual é o impacto real
do tempo de startup com compilação dinâmica e estática? Na HotSpot, quem
decide entre interpretador e JIT é a política de compilação, descrita no
comentário de cabeçalho de \citetitle{doc:hotspot-compilationpolicy}. Ela
tem cinco níveis: o 0 é o interpretador, o 1, o 2 e o 3 são o C1 com
graus diferentes de coleta de perfil, e o 4 é o C2, \enquote{with full
profile guided optimization}. A execução começa no nível~0, e a política
pode compilar o método no nível~3 ou no~2; terminada a coleta de perfil
no nível~3, o método passa ao nível~4. A passagem do nível~0 ao~3 compara
o número de invocações e de \emph{back-edges} (voltas de laço) do método
com limiares multiplicados por um coeficiente de escala, e o tamanho das
filas do C1 e do C2 ajusta esses limiares. Experimento (JDK~25.0.4.1,
imagem \texttt{eclipse-temurin:25-jdk}): \texttt{java
-XX:+PrintFlagsFinal -version} mostrou \texttt{Tier3InvocationThreshold}
igual a 200 e \texttt{Tier4InvocationThreshold} igual a 5000.

O AOT aparece de dois jeitos. Dentro da HotSpot, o
\citetitle{doc:jep483} (JDK~24) guarda num cache as classes carregadas e
ligadas (seção~\ref{sec:class-loading}), e o \citetitle{doc:jep515}
(JDK~25) guarda também perfis de métodos coletados num treino, para que o
JIT gere código nativo logo no início. Fora dela, o
\citetitle{doc:graalvm-native-image} compila o código Java, antes da
execução, para um executável nativo (seção~\ref{sec:jvm-vs-nativo}).

Os números publicados vêm dos próprios JEPs. O JEP~483 informa que um
programa pequeno roda em 0,031~s no JDK~23 e em 0,018~s no JDK~24 com o
cache, 42\% mais rápido; o número do Spring PetClinic está na
seção~\ref{sec:class-loading}. O JEP~515 informa que o seu programa de
exemplo roda em 90~ms com um cache sem perfis e em 73~ms depois de
coletar os perfis no cache, 19\% mais rápido. Para JVM contra executável
nativo, a seção~\ref{sec:jvm-vs-nativo} traz a tabela que o Quarkus
publica, com tempo até a primeira requisição, throughput de pico e
memória nos dois modos e com o cache AOT.

\paragraph{Otimizações que dependem do perfil.}
\begin{itemize}
\item \emph{Inlining}: o corpo do método chamado é copiado no chamador.
  É a porta para as outras otimizações, porque o compilador passa a
  enxergar os dois métodos juntos.
\item \termo{Escape analysis}{Análise que descobre se um objeto
  \enquote{escapa} do método ou da thread que o criou.}: se o objeto não
  escapa, o C2 pode eliminar a alocação (\emph{scalar replacement}) e os
  locks sobre ele (\emph{lock elision}). A documentação da Oracle
  ressalta que a HotSpot \emph{não} troca alocação no heap por alocação
  na pilha (\citetitle{doc:hotspot-perf}).
\item Especulação: o C2 compila só os caminhos que o perfil mostrou
  vivos e deixa nos outros uma \termo{uncommon trap}{Ponto no código
  compilado que devolve o controle ao runtime quando um caminho
  considerado impossível acontece.}. Se o caminho raro acontece, o
  código compilado é descartado
  (\termo{deoptimization}{Desotimização: trocar, em plena execução,
  frames de código otimizado por frames equivalentes de código não
  otimizado.}), o método volta
  ao interpretador ou a uma camada inferior, ganha novo perfil e é
  recompilado (\citetitle{artigo:shipilev-uncommon-traps}).
\item Chamadas virtuais: num ponto de chamada que só viu um tipo, o
  compilador pode \emph{devirtualizar}; com muitos tipos (megamórfico),
  a chamada vai por um despacho mais lento. Shipilëv mede a diferença e
  conclui que ela não justifica evitar chamadas virtuais
  (\citetitle{artigo:shipilev-megamorphic}).
\end{itemize}

\paragraph{Safepoints.}
Várias operações da JVM (algumas fases do GC, desotimização) precisam
parar as threads Java em pontos conhecidos, os
\termo{safepoints}{Pontos do código em que o estado da thread é
conhecido pela JVM; as threads verificam periodicamente se devem parar
ali.}. O código compilado tem verificações baratas (\emph{safepoint
polls}) que quase nunca disparam
(\citetitle{artigo:shipilev-safepoint}). O artigo explica por que os
polls existem: uma thread num laço ocupado pode nunca chamar a JVM por
conta própria, e por isso a HotSpot usa um esquema cooperativo, em que a
thread pergunta de tempos em tempos se deve parar num ponto de estado
conhecido. A opção \texttt{-Xlog:safepoint} registra os
safepoints no log da JVM.

\paragraph{As decisões que isso muda.}
\begin{itemize}
\item \emph{Warm-up é real.} As primeiras requisições rodam interpretadas
  ou em C1. Shipilëv avisa que o aquecimento precisa ter o
  mesmo perfil da produção (\citetitle{artigo:shipilev-uncommon-traps});
  senão, a recompilação aparece depois, no
  meio da medição. Teste de carga sem aquecimento mede o interpretador.
\item \emph{Microbenchmark à mão engana.} Um \texttt{System.nanoTime()}
  em volta de um laço mede compilação, eliminação de código morto e
  especulação. Use o \citetitle{doc:jmh}, cuja página pede que os
  benchmarks sejam revisados por pares.
\item \emph{Escreva código claro primeiro.} Métodos pequenos e objetos
  temporários que não escapam são exatamente o que o C2 otimiza melhor;
  o item 67 do \citetitle{livro:bloch2017} (\enquote{Optimize
  judiciously}) chega à mesma conclusão por outro caminho.
\end{itemize}

\paragraph{Vendo o JIT trabalhar.}
O manual do \texttt{java} documenta \texttt{-XX:+PrintCompilation}
(imprime uma linha a cada método compilado), \texttt{-Xint} (só
interpretador, para comparar), \texttt{-XX:-TieredCompilation} e, após
\texttt{-XX:+UnlockDiagnosticVMOptions}, \texttt{-XX:+PrintInlining} e
\texttt{-XX:+LogCompilation} (log em \texttt{hotspot\_pid<pid>.log} no
diretório atual; o manual do \texttt{java} ainda diz \texttt{hotspot.log},
e \texttt{-XX:LogFile} escolhe outro nome), que o
\citetitle{doc:jitwatch} visualiza). A seção~\ref{sec:ferramentas}
reúne essas ferramentas.
